% !TeX root = elegantnote-en.tex
\section{Carboxylic Acid Derivatives}

\subsection{Introduction 20260519}

\begin{enumerate}
    \item The most important acid derivatives are esters, amides, and nitriles. Acid halides and anhydrides are also included.
\end{enumerate}

\subsection{Structure and Nomenclature of Acid Derivatives 20260519 20260521}

\subsubsection{Esters of Carboxylic Acids}

\begin{enumerate}
    \item Esters are carboxylic acid derivatives in which the hydroxyl group (\chemfig{-OH}) is replaced by an alkoxy group (\chemfig{-OR}).
    \item Esters can be formed by the Fischer esterification of an acid with an alcohol.
    \item The names of esters consist of two words that reflect their composite structure. The first word is derived from the \textit{alkyl} group of the alcohol, and the second word from the \textit{carboxylate} group of the carboxylic acid.
    \item Lactones
    \begin{enumerate}
        \item Cyclic esters are called \textbf{lactones}.A lactone is formed from an openchain hydroxy acid in which the hydroxyl group has reacted with the acid group to form an ester.
        \item The IUPAC names of lactones are derived by adding the term \textit{lactone} at the end of the name of the parent acid. (open the ring and name it.)
        \item The common names of lactones are formed by changing the \textit{-ic} acid ending of the hydroxy acid to \textit{-olactone}.

    \end{enumerate}
\end{enumerate}

\subsubsection{Amides}

\begin{enumerate}
    \item An amide is a composite of a carboxylic acid and ammonia or an amine.
    \item Amides are only weakly basic, and we consider the amide functional group to be neutral.
    \item \textbf{primary amide}, \textbf{secondary amide} or \textbf{an N-substituted amide}, \textbf{tertiary amides} or \textbf{N,N-disubstituted amides}.
    \item Nomenclature: Drop the \textit{-ic} acid or \textit{-oic} acid suffix, and add the suffix \textit{-amide}. For acids that are named as alkanecarboxylic acids, the amides are named by using the suffix \textit{-carboxamide}.
    \item Lactams
    \begin{enumerate}
        \item Cyclic amides are called \textbf{lactams}.
        \item Lactams are named like lactones, by adding the term \textit{lactam} at the end of the IUPAC name of the parent acid.
        \item Common names of lactams are formed by changing the \textit{-ic} acid ending of the amino acid to \textit{-olactam}.
    \end{enumerate}
\end{enumerate}

\subsubsection{Nitriles}

\begin{enumerate}
    \item \textbf{Nitriles} contain the \textbf{cyano group}, \chemfig{-C~N}.
    \item They hydrolyze to give carboxylic acids and can be synthesized by dehydration of amides.
    \item Hydrolysis to an acid (\ce{H2O}), Synthesis from an acid (\ce{POCl3}).
    \item Both the carbon atom and the nitrogen atom of the cyano group are \textit{sp} hybridized, and the \chemfig{R-C~N} bond angle is \qty{180}{\degree} (linear). The structure is similar to that of a terminal alkyne, except that the nitrogen atom  has a lone pair of electrons.
    \item Common names of nitriles begin with the common name of the acid, and replace the suffix \textit{-ic acid} with the suffix \textit{-onitrile}. The IUPAC name is constructed from the alkane name, with the suffix \textit{-nitrile} added.
\end{enumerate}

\subsubsection{Acid Halides}

\begin{enumerate}
    \item The most common acyl halides are the acid chlorides (acyl chlorides).
    \item An acid halide is named by replacing the \textit{-ic acid} suffix of the acid name (either the common name or the IUPAC name) with \textit{-yl} and the halide name. For acids that are named as alkanecarboxylic acids, the acid chlorides are named by using the suffix \textit{-carbonyl chloride}.
\end{enumerate}

\subsubsection{Acid Anhydrides}

\begin{enumerate}
    \item An acid anhydride contains two molecules of an acid, with loss of a molecule of water.
    \item Anhydride nomenclature: the word \textit{acid} is changed to \textit{anhydride} in both the common name and the IUPAC name.
    \item Anhydrides composed of two different acids are called \textit{mixed anhydrides} and are named by using the names of the individual acids.
\end{enumerate}

\subsubsection{Nomenclature of Multifunctional Compounds}

\begin{enumerate}
    \item acid > ester > amide > nitrile > aldehyde > ketone > alcohol > amine > alkene, alkyne
    \item ethyl \textit{o}-cyanobenzoate, 2-formylcyclohexanecarboxamide, 2-hydroxybutanenitrile
    \item Summary of Functional Group Nomenclature
\end{enumerate}

\begin{example} % 21-1 !!!
    Name the following carboxylic acid derivatives, giving both a common name and an IUPAC name where possible.
\end{example}

\subsection{Physical Properties of Carboxylic Acid  Derivatives 20260521}

\subsubsection{Boiling Points and Melting Points}

\begin{enumerate}
    \item 
\end{enumerate}

\subsubsection{Solubility}

\subsection{Spectroscopy of Carboxylic Acid Derivatives 20260521}

\subsubsection{Infrared Spectroscopy}

\subsubsection{NMR Spectroscopy}

\begin{example} % 21-5
    For each set of IR and NMR spectra, determine the structure of the unknown compound. Explain how your proposed structure fits the spectra.
\end{example}

\subsection{Interconversion of Acid Derivatives by Nucleophilic Acyl Substitution 20260521 20260526}

\begin{mechanism}
    Addition–Elimination Mechanism of Nucleophilic Acyl Substitution
\end{mechanism}

\subsubsection{Reactivity of Acid Derivative}

\begin{enumerate}
    \item Acid derivatives differ greatly in their reactivity toward nucleophilic acyl substitution.
    \item The order of reactivity stems partly from the basicity of the leaving groups. %pic !!!
\end{enumerate}

\subsubsection{Favorable Interconversions of Acid Derivatives}

\begin{enumerate}
    \item More reactive acid derivatives are easily converted to less reactive derivatives. %pic !!!
\end{enumerate}

\begin{mechanism}
    Conversion of an Acid Chloride to an Anhydride

    This mechanism follows the standard pattern of an addition–elimination mechanism, ending with loss of a proton to give the final product.
    \begin{itemize}
        \item Step 1: Addition of the nucleophile.
        \item Step 2: Elimination of the leaving group.
        \item Step 3: Loss of a proton.
    \end{itemize}
\end{mechanism}

\begin{mechanism}
    Conversion of an Acid Chloride to an Ester

    This is another reaction that follows the standard addition–elimination mechanism, ending with loss of a proton to give the final product.
\end{mechanism}

\begin{mechanism}
    Conversion of an Acid Chloride to an Amide

    This reaction also follows the steps of a standard addition–elimination mechanism, ending with loss of a proton to give the amide.
\end{mechanism}

\begin{mechanism}
    Conversion of an Acid Anhydride to an Ester
    
    This reaction follows the standard addition–elimination mechanism, ending with loss of a proton to give the ester.
\end{mechanism}

\begin{mechanism}
    Conversion of an Acid Anhydride to an Amide

    This reaction follows the standard addition–elimination mechanism, ending with loss of a proton to give the amide.
\end{mechanism}

\begin{mechanism}
    Conversion of an Ester to an Amide (Ammonolysis of an Ester)

    This is yet another standard addition–elimination mechanism, ending with loss of a proton to give the amide.
\end{mechanism}

\subsubsection{Leaving Groups in Nucleophilic Acyl Substitutions}

\begin{enumerate}
    \item The reaction’s one-step mechanism is not strongly endothermic or exothermic.
    \item In the acyl substitution, the leaving group leaves in a separate second step. This second step is highly exothermic. 
    \item A strong base may serve as a leaving group if it leaves in a highly exothermic step, usually converting an unstable, negatively charged intermediate to a stable molecule.
\end{enumerate}

\begin{example} % 21-7
    Which of the following proposed reactions would take place quickly under mild conditions?
\end{example}

\subsection{Transesterification 20260526}

\begin{enumerate}
    \item Esters undergo \textbf{transesterification}, in which one alkoxy group substitutes for another, under either acidic or basic conditions.
    \item An equilibrium results, and the equilibrium can be driven toward the desired ester by using a large excess of the desired alcohol or by removing the other alcohol.
    \item Base-Catalyzed Transesterification
    \item Acid-Catalyzed Transesterification
    \begin{enumerate}
        \item Consider the carbon skeletons of the reactants and products, and identify which  carbon atoms in the products are likely derived from which carbon atoms in the reactants.
        \item Consider whether any of the reactants is a strong enough electrophile to react without  being activated. If not, consider how one of the reactants might be converted to a strong electrophile by protonation of a Lewis basic site.
        \item Consider how a nucleophilic site on another reactant can attack the strong electrophile  to form a bond needed in the product.
        \item Consider how the product of nucleophilic attack might be converted to the final product  or reactivated to form another bond needed in the product.
        \item Draw out all steps of the mechanism, using curved arrows to show the movement  of electrons.
    \end{enumerate}
\end{enumerate}

\begin{example} % 21-12
    When ethyl 4-hydroxybutyrate is heated in the presence of a trace of a basic catalyst (sodium acetate), one of the products is a lactone. Propose a mechanism for formation of this lactone.
\end{example}

\begin{mechanism} 
    Transesterification

    Following is a summary of the mechanism of transesterification under basic and acidic conditions. (Aspiring)
    \begin{enumerate}
        \item Base-catalyzed
        
        Base-catalyzed transesterification is a simple two-step nucleophilic acyl substitution。
        \item Acid-catalyzed
        
        Acid-catalyzed transesterification requires extra proton transfers before and after the major steps. The overall reaction takes place in two stages. The first half of the reaction involves acid-catalyzed addition of the nucleophile, and the second half involves acid-catalyzed elimination of the leaving group.
    \end{enumerate}
\end{mechanism}

\subsection{Hydrolysis of Carboxylic Acid Derivatives 20260526}

\subsubsection{Hydrolysis of Acid Halides and Anhydrides}

\begin{enumerate}
    \item Acid halides and anhydrides are so reactive that they hydrolyze under neutral conditions.
\end{enumerate}

\subsubsection{Hydrolysis of Esters}

\begin{enumerate}
    \item Acid-catalyzed hydrolysis of an ester is simply the reverse of the Fischer esterification equilibrium.
    \item Basic hydrolysis of esters, called \textbf{saponification}, avoids the equilibrium of the Fischer esterification.
\end{enumerate}

\begin{mechanism}
    Saponification of an Ester

    This is another standard addition–elimination mechanism, ending with a proton transfer to give the final products.
\end{mechanism}

\begin{example} % 21-16
    Suppose we have some optically pure (\textit{R})-2-butyl acetate that has been “labeled” with the heavy \ce{^18O} isotope at one oxygen atom as shown.
\end{example}

\subsubsection{Hydrolysis of Amides}

\begin{enumerate}
    \item 
\end{enumerate}

\begin{mechanism}
    Basic Hydrolysis of an Amide
\end{mechanism}

\begin{mechanism}
    Acidic Hydrolysis of an Amide

    This mechanism takes place in two stages.
\end{mechanism}

\subsubsection{Hydrolysis of Nitriles}

\begin{enumerate}
    \item Nitriles are hydrolyzed to amides, and further to carboxylic acids, by heating with aqueous acid or base.
    \item Mild conditions can hydrolyze a nitrile only as far as the amide. Stronger conditions can hydrolyze it all the way to the carboxylic acid.
\end{enumerate}

\begin{mechanism}
    Base-Catalyzed Hydrolysis of a Nitrile


\end{mechanism}

\begin{example} % 21-23
    The mechanism for acidic hydrolysis of a nitrile resembles the basic hydrolysis, except that the nitrile is first protonated, activating it toward attack by a weak nucleophile (water). Under acidic conditions, the proton transfer (tautomerism) involves protonation on nitrogen followed by deprotonation on oxygen. Propose a mechanism for the acid-catalyzed hydrolysis of benzonitrile to benzamide.
\end{example}

\subsection{Reduction of Acid  Derivatives 20260526}

\subsubsection{Reduction to Alcohols}

\begin{enumerate}
    \item Lithium aluminum hydride reduces acids, acid chlorides, anhydrides, and esters to primary alcohols.
\end{enumerate}

\begin{mechanism}
    Hydride Reduction of an Ester
\end{mechanism}

\subsubsection{Reduction to Aldehydes}

\subsubsection{Reduction to Amines}

\begin{enumerate}
    \item % 亚甲基而不是醇
\end{enumerate}

\subsection{Reactions of Acid Derivatives with Organometallic Reagents 20260526}

\begin{enumerate}
    \item Esters and Acid Chlorides
    \begin{enumerate}
        \item Grignard and organolithium reagents add twice to acid chlorides and esters to give alkoxides (Section 10-9D). Protonation of the alkoxides gives alcohols. % Grignard not react with acid.
    \end{enumerate}
    \item Nitriles
    \begin{enumerate}
        \item A Grignard or organolithium reagent attacks the electrophilic cyano group to form the salt of an imine.
        \item Acidic hydrolysis of the salt (in a separate step) gives the imine, which is further hydrolyzed to a ketone (Section 18-9).
    \end{enumerate}
\end{enumerate}

\begin{example} % 21-26
    Acidic hydrolysis of the salt (in a separate step) gives the imine, which is further hydrolyzed to a ketone (Section 18-9).
\end{example}

\subsection{Summary of the Chemistry of Acid Chlorides 20260526}

\begin{enumerate}
    \item Synthesis of Acid Chlorides: Acid chlorides (acyl chlorides) are synthesized from the corresponding carboxylic acids, \ce{SOCl2} and \ce{COCl2} most convenient.
    \item Reactions of Acid Chlorides
    \item Friedel–Crafts Acylation of Aromatic Rings
\end{enumerate}

\subsection{Summary of the  Chemistry of Anhydrides 20260526}

\begin{enumerate}
    \item l
    \item Acetic Anhydride: The most common industrial synthesis begins by dehydrating acetic acid to give ketene.
    \item General Anhydride Synthesis: The most general method for making anhydrides is the reaction of an acid chloride with a carboxylic acid or a carboxylate salt.
    \item Reactions of Anhydrides: Like acid chlorides, anhydrides participate in the Friedel–Crafts acylation. Cyclic anhydrides can provide additional functionality on the side chain of the aromatic product.
    \item Three specific instances when anhydrides are preferred.
    \begin{enumerate}
        \item 
    \end{enumerate}
\end{enumerate}

\begin{example}
    \begin{enumerate}[label=(\alph*)]
        \item Give the products expected when acetic formic anhydride reacts with (i) aniline and  (ii) benzyl alcohol. 
        \item Propose mechanisms for these reactions.
    \end{enumerate}
\end{example}

\subsection{Summary of the Chemistry of Esters 20260526}

\begin{enumerate}
    \item Synthesis of Esters
    \item Reactions of Esters
    \item Formation of Lactones
\end{enumerate}

\subsection{Summary of the Chemistry of Amides 20260526}

\begin{enumerate}
    \item Synthesis of Amides
    \item Reactions of Amides
    \item Dehydration of Amides to Nitriles: Strong dehydrating agents can remove the elements of water from a primary amide to give a nitrile.
    \item Formation of Lactams
    \item Biological Reactivity of $\beta$-Lactams: penicillin
    \item Polyamides: Nylon
\end{enumerate}

\subsection{Summary of the Chemistry of Nitriles 20260528}

\begin{enumerate}
    \item Synthesis
    \item Reactions of Nitriles
\end{enumerate}

\subsection{Thioesters 20260528}

\begin{enumerate}
    \item A \textbf{thioester} is formed from a carboxylic acid and a thiol.
    \item Thioesters are more reactive toward nucleophilic acyl substitution than normal esters, but less reactive than acid chlorides and anhydrides.
    \item The overlap is weak and relatively ineffective.
\end{enumerate}

\subsection{Esters and Amides of Carbonic Acid 20260528}

\begin{enumerate}
    \item \textbf{Carbonic acid} is formed reversibly whenever carbon dioxide dissolves in water.
    \item \textbf{Carbonate esters} are diesters of carbonic acid, with two alkoxy groups replacing the hydroxyl groups of carbonic acid.
    \item \textbf{Ureas} are diamides of carbonic acid, with two nitrogen atoms bonded to the carbonyl group.
    \item \textbf{Carbamate esters} (\textbf{urethanes}) are the stable esters of the unstable \textbf{carbamic acid}, the monoamide of carbonic acid.
    \item Many of these derivatives can be synthesized by nucleophilic acyl substitution from phosgene, the acid chloride of carbonic acid. Another way of making urethanes is to treat an alcohol or a phenol with an \textbf{isocyanate}, which is an anhydride of a carbamic acid.
    \item Polycarbonates and Polyurethanes
    \begin{enumerate}
        \item two large classes of polymers are bonded by linkages containing these functional groups: the \textbf{polycarbonates} and the \textbf{polyurethanes}.
        \item \textit{bisphenol A}
    \end{enumerate}
\end{enumerate}

%44 47 49/bc 50

\begin{example} % 21-64
    An unknown compound gives a mass spectrum with a weak molecular ion at \textit{m/z} 113 and a prominent ion at \textit{m/z} 68. Its  NMR and IR spectra are shown here. Determine the structure, and show how it is consistent with the observed absorptions. Propose a favorable fragmentation to explain the prominent MS peak at \textit{m/z} 68.

    \textbf{SOLUTION} \quad 
    
    \textbf{MS} 113-odd-\ce{N} 113 - 68 = 45 

    \textbf{IR} 3000 medium \ce{-NH2}?, 2200 triple bond, 1750 carbonyl, 1200 \ce{C-O}

    \textbf{C NMR} 5 \ce{C}, 5 * 12 = 60

    till now, C5 N O2 sum = 106, so 7H

    \textbf{H NMR} one single \ce{-CH2}, one \ce{-CH2} and \ce{-CH3} together (split).

    \chemfig{N~C-CH_2-C(=[2]O)-O-CH_2-CH_3} \ce{-OEt} = 45
\end{example}

\begin{example} % 21-65
    An unknown compound gives the NMR, IR, and mass spectra shown next. Propose a structure, and show how it is  consistent with the observed absorptions. Show fragmentations that account for the prominent ion at \textbf{m/z} 69 and the smaller peak at \textbf{m/z} 99.

    \textbf{SOLUTION} \quad 

    \textbf{MS} 114 - 69 = 45(\ce{-OEt}), 114 - 99 = 13(\ce{-CH3}), no benzyl(91)

    \textbf{IR} 3000 no \ce{-OH} \ce{-NH2} 1700-1600 \chemfig{C=C-C(=[2]O)} classic

    \textbf{C NMR} 5 \ce{C}? 6?(because of 5 kind of H)
    
    \textbf{H NMR} 6-7 H of \ce{C=C} beside it is \ce{-CH3}, the rest is \ce{-OEt}.
    
    notice: E not Z % 耦合常数显示为反式构象

    \chemfig{C(-[3]H)(-[5]H_3C)=C(-[1]C(=[2]O)-OC_2H_5)(-[7]H)}
\end{example}

\begin{example} % 21-66
    The IR spectrum, \ce{^13C} NMR spectrum, and \ce{^1H} NMR spectrum of an unknown compound (\ce{C6H8O3}) appear next.  Determine the structure, and show how it is consistent with the spectra.

    \textbf{SOLUTION} \quad

    $\Omega$ = 3, 2 carbonyl and 1 ring? no anhydride(1800)

    \textbf{C NMR} one over 200
    
    \textbf{H NMR} 3:1:1:1:2 the highest one not splitted, \ce{-CH3} to \ce{C=O}.
\end{example}

\begin{example} % 21-67
    An unknown compound of molecular formula \ce{C5H9NO} gives the IR and NMR spectra shown here. The broad NMR peak at $\delta 7.55$ disappears when the sample is shaken with \ce{D2O}. Propose a structure, and show how it is consistent with the absorptions in the spectra.
\end{example}